% Original, candidate11 and proposed complete proofs; checking archive only.

% inv:B3.1 ORIGINAL
\begin{proposition}[Connected cyclic covers and the cusp comparison]\hypertarget{tag-01S7}{}\label{inv:B3.1}
\StatementRecord{inv:B3.1}{proved}{01S7}{\SeriesRef{inv:B3.3}, \SeriesRef{inv:B1.1}}

For an admissible $(3,4)$ member with
$n=|12l_0-4l_1-3l_2|>1$, let $\pi:Y\to X$ be the universal
cover and $g=f\pi$. It is a finite etale holomorphic map of degree $n$.
Every fibre of $g$ is connected; its finite reductions are bielliptic,
its multiplicities are $3,4$, and its cusp $W_n$ normalizes to a
disjoint union of $n$ copies of $\mathrm{dP}_6$.
The pullback maps induce isomorphisms
\[
R^qf_*\OO_X\simeq R^qg_*\OO_Y,\qquad
R^qf_*E_X\simeq R^qg_*E_Y.
\]
\end{proposition}


\begin{proof}
\begin{enumerate}[label={$\langle1\rangle$\arabic*.},leftmargin=*]
\item We first derive the fundamental-group presentation from the fillings.
The smooth family over the punctured base has zero section, so its
fundamental group is $\Lambda\rtimes\langle x,y\rangle$, where
$x,y$ are the analytic peripheral generators fixed by the root-deck convention and act by $A_1,A_2$.
On a finite filling the lifted affine generator has $m_j$th power
translation by $v_j$: its $m_j$ translations add to $v_j$ because
$A_jv_j=v_j$. Thus filling the normal disc imposes
$x^3=t_{v_1}$ and $y^4=t_{v_2}$, with the signs fixed by the
clockwise convention on the first page.

\item The toric covering space of the cusp tube is simply connected.
Indeed its dense torus has cocharacter lattice $\Z^3$; a loop may
be perturbed into that torus, and the normal discs of the three rays
$(0,1),(e_1,1),(e_2,1)$ kill a basis of this lattice.
This argument applies also to the infinite fan, since a compact loop
and a homotopy meet only finitely many of its locally finite charts.
The same conclusion holds for the small tube: if a loop in
$|t_c|<\epsilon$ bounds a disc in the toric space, the torus action
$(x_1,x_2,t_c)\mapsto(x_1,x_2,a t_c)$ for a sufficiently small
positive $a$ moves that compact disc into the tube, and the scaling
homotopy of its boundary stays in the tube.
The cusp quotient therefore has fundamental group
$\Lambda/\Lambda_{\rm tor}$, and its toric meridian bounds the
normal disc at fixed $x_1,x_2$. A translated collar makes this
meridian $t_\mu y^{-1}x^{-1}$, where $\gamma(\mu)=l_0$.
Van Kampen consequently kills $\Lambda_{\rm tor}$ and imposes
$xy=t_\mu$.

\item The normal closure of $\Lambda_{\rm tor}$ in the fibre lattice is
$\ker\gamma$: it contains $\hat w,\hat\delta$ and
$A_1\hat w+\hat w=\hat u$, and all these relations lie in
$\ker\gamma$. The remaining generator $c$, the image of
$\hat\gamma$, is central because $\gamma A_j=\gamma$.
The presentation is thus
$\langle c,x,y\mid c\text{ central},\ xy=c^{l_0},\
x^3=c^{l_1},\ y^4=c^{l_2}\rangle$.
Eliminating $y$ makes it abelian, with relation matrix
\[
 \begin{pmatrix}3&-l_1\\4&l_2-4l_0\end{pmatrix},
 \qquad\det=-p.
\]
Its first Smith divisor is one because its entries include three
and four; hence the group is $\Z/|p|$.
After imposing $c=1$ the quotient has $x^3=x^4=1$, so it is trivial.
Thus the image of the regular-fibre lattice, containing $c$, generates
the entire group $\Z/n$. The same is true of each reduced special
fibre. At a finite point the cyclic model is a smooth torus family
over a disc, with a free affine finite-group action. An equivariant
smooth trivialization, obtained by averaging a horizontal connection,
contracts the disc to its centre and descends to a deformation
retraction of the local total space onto $S_j$. Thus the image of
the regular-fibre lattice in the local fundamental group is contained
in the image of $\pi_1(S_j)$; the same holds after mapping to
$\pi_1(X)$. At the cusp $\hat\gamma$ survives in
$\Lambda/\Lambda_{\rm tor}$. The explicit normal-crossing contraction
of the tube onto $W$ given in \ref{inv:B3.3} identifies their fundamental groups,
so the same generator is in the image of $\pi_1(W)$.
The covering-space criterion therefore makes the inverse image of
every fibre connected. Admissibility gives $3\nmid l_1$ and odd $l_2$,
so $\gcd(n,12)=1$. The etale cover preserves fibre multiplicities.
On a finite reduction it cannot kill the nontrivial order-$3$ or
order-$4$ linear holonomy, since the covering degree is prime to it;
the resulting finite torus quotient is again bielliptic.

\item Normalization commutes with the finite etale base change $Y\to X$.
The pullback of the simply connected normalization $\mathrm{dP}_6$
is its trivial degree-$n$ cover. The double curves and triple points
also pull back etale, and their incidence maps are the lifts of the
original incidence maps. Consequently the dual complex is an actual
degree-$n$ covering of the original real two-torus, rather than merely
a cell complex with the same cell counts. It is connected because
$W_n$ is connected. The normalization resolution for $\OO_{W_n}$
has acyclic components, double curves and point strata. Its global
cochain complex is consequently the cellular cochain complex of this
dual torus, with $n,3n,2n$ cells. Thus its coherent cohomology is
$(1,2,1)$, and pullback from $W$ is an isomorphism over $\C$:
on a finite cover of real tori it is induced by a full-rank map of
their integral first cohomology lattices.

\item Modulo $d\log q$, the two toric forms $d\log x_i$ are invariant
under all cusp gluings as computed in \ref{inv:B1.1}. Hence
$E_X|_W=\OO_W^2$ and $E_Y|_{W_n}=\OO_{W_n}^2$,
and the same pullback argument applies to their cohomology.
Over a regular fibre, an isogeny of complex tori induces isomorphisms
on $H^q(\OO)$ and $H^q(\Omega^1)$ over $\C$.
At a finite fibre, make the cyclic base change. The map is an isogeny
of the smooth torus families and induces isomorphisms of their
cohomology bundles. Taking the same cyclic invariant eigenspaces
descends these isomorphisms, including all nilpotent layers of the
scheme fibre. Grauert base change at the cusp, and this explicit
calculation at the finite points, prove the asserted bundle maps are
isomorphisms on all of $B$.
\end{enumerate}
\end{proof}

% CANDIDATE11
\begin{proposition}[Connected cyclic covers and the cusp comparison]\hypertarget{tag-01S7}{}\label{inv:B3.1}
\StatementRecord{inv:B3.1}{proved}{01S7}{\SeriesRef{inv:B3.3}, \SeriesRef{inv:B1.1}}

For an admissible $(3,4)$ member with
$n=|12l_0-4l_1-3l_2|>1$, let $\pi:Y\to X$ be the universal
cover and $g=f\pi$. It is a finite \'etale holomorphic map of degree $n$.
Every fibre of $g$ is connected; its finite reductions are bielliptic,
its multiplicities are $3,4$, and its cusp $W_n$ normalizes to a
disjoint union of $n$ copies of $\mathrm{dP}_6$.
For $\mathcal E_X$ of Definition~\ref{I:def:coarse-log-cotangent}
and the identical quotient $\mathcal E_Y$ on $Y$, pullback induces isomorphisms
\[
R^qf_*\OO_X\simeq R^qg_*\OO_Y,\qquad
R^qf_*\mathcal E_X\simeq R^qg_*\mathcal E_Y.
\]
\end{proposition}


\begin{proof}
\begin{enumerate}[label=\arabic*.,leftmargin=*]
\item We first derive the fundamental-group presentation from the fillings.
The smooth family over the punctured base has zero section, so its
fundamental group is $\Lambda\rtimes\langle x,y\rangle$, where
$x,y$ are the analytic peripheral generators fixed by the root-deck convention and act by $A_1,A_2$.
On a finite filling the lifted affine generator has $m_j$th power
translation by $v_j$: its $m_j$ translations add to $v_j$ because
$A_jv_j=v_j$. Thus filling the normal disc imposes
$x^3=t_{v_1}$ and $y^4=t_{v_2}$, with the signs fixed by the
clockwise convention on the first page.

\item The toric covering space of the cusp tube is simply connected.
Indeed its dense torus has cocharacter lattice $\Z^3$; a loop may
be perturbed into that torus, and the normal discs of the three rays
$(0,1),(e_1,1),(e_2,1)$ kill a basis of this lattice.
This argument applies also to the infinite fan, since a compact loop
and a homotopy meet only finitely many of its locally finite charts.
The same conclusion holds for the small tube: if a loop in
$|t_c|<\epsilon$ bounds a disc in the toric space, the torus action
$(x_1,x_2,t_c)\mapsto(x_1,x_2,a t_c)$ for a sufficiently small
positive $a$ moves that compact disc into the tube, and the scaling
homotopy of its boundary stays in the tube.
The cusp quotient therefore has fundamental group
$\Lambda/\Lambda_{\rm tor}$, and its toric meridian bounds the
normal disc at fixed $x_1,x_2$. A translated collar makes this
meridian $t_\mu y^{-1}x^{-1}$, where $\gamma(\mu)=l_0$.
Van Kampen consequently kills $\Lambda_{\rm tor}$ and imposes
$xy=t_\mu$.

\item The normal closure of $\Lambda_{\rm tor}$ in the fibre lattice is
$\ker\gamma$: it contains $\hat w,\hat\delta$ and
$A_1\hat w+\hat w=\hat u$, and all these relations lie in
$\ker\gamma$. The remaining generator $c$, the image of
$\hat\gamma$, is central because $\gamma A_j=\gamma$.
The presentation is thus
$\langle c,x,y\mid c\text{ central},\ xy=c^{l_0},\
x^3=c^{l_1},\ y^4=c^{l_2}\rangle$.
Eliminating $y$ makes it abelian, with relation matrix
\[
 \begin{pmatrix}3&-l_1\\4&l_2-4l_0\end{pmatrix},
 \qquad\det=-p.
\]
Its first Smith divisor is one because its entries include three
and four; hence the group is $\Z/|p|$.
After imposing $c=1$ the quotient has $x^3=x^4=1$, so it is trivial.
Thus the image of the regular-fibre lattice, containing $c$, generates
the entire group $\Z/n$. The same is true of each reduced special
fibre. At a finite point the cyclic model is a smooth torus family
over a disc, with a free affine finite-group action. An equivariant
smooth trivialization, obtained by averaging a horizontal connection,
contracts the disc to its centre and descends to a deformation
retraction of the local total space onto $S_j$. Thus the image of
the regular-fibre lattice in the local fundamental group is contained
in the image of $\pi_1(S_j)$; the same holds after mapping to
$\pi_1(X)$. At the cusp $\hat\gamma$ survives in
$\Lambda/\Lambda_{\rm tor}$. The explicit normal-crossing contraction
of the tube onto $W$ given in \ref{inv:B3.3} identifies their fundamental groups,
so the same generator is in the image of $\pi_1(W)$.
The covering-space criterion therefore makes the inverse image of
every fibre connected. Admissibility gives $3\nmid l_1$ and odd $l_2$,
so $\gcd(n,12)=1$. The etale cover preserves fibre multiplicities.
On a finite reduction it cannot kill the nontrivial order-$3$ or
order-$4$ linear holonomy, since the covering degree is prime to it;
the resulting finite torus quotient is again bielliptic.

\item Normalization commutes with the finite etale base change $Y\to X$.
The pullback of the simply connected normalization $\mathrm{dP}_6$
is its trivial degree-$n$ cover. The double curves and triple points
also pull back etale, and their incidence maps are the lifts of the
original incidence maps. The lifted incidence maps therefore make the dual complex a
degree-$n$ covering of the original real two-torus. It is connected because
$W_n$ is connected. The normalization resolution for $\OO_{W_n}$
has acyclic components, double curves and point strata. Its global
cochain complex is consequently the cellular cochain complex of this
dual torus, with $n,3n,2n$ cells. Thus its coherent cohomology is
$(1,2,1)$, and pullback from $W$ is an isomorphism over $\C$:
on a finite cover of real tori it is induced by a full-rank map of
their integral first cohomology lattices.

\item Modulo $d\log q$, the two toric forms $d\log x_i$ are invariant
under all cusp gluings as computed in \ref{inv:B1.1}. Hence
$\mathcal E_X|_W=\OO_W^2$ and $\mathcal E_Y|_{W_n}=\OO_{W_n}^2$,
and the same pullback argument applies to their cohomology.
Over a regular fibre, an isogeny of complex tori induces isomorphisms
on $H^q(\OO)$ and $H^q(\Omega^1)$ over $\C$.
At a finite fibre, make the cyclic base change. The map is an isogeny
of the smooth torus families and induces isomorphisms of their
cohomology bundles. Taking the same cyclic invariant eigenspaces
descends these isomorphisms, including all nilpotent layers of the
scheme fibre. Grauert base change at the cusp, and this explicit
calculation at the finite points, prove the asserted bundle maps are
isomorphisms on all of $B$.
\end{enumerate}
\end{proof}

% PROPOSED
\begin{proposition}[Connected cyclic covers and the cusp comparison]\hypertarget{tag-01S7}{}\label{inv:B3.1}
\StatementRecord{inv:B3.1}{proved}{01S7}{\cref{I:lem:cusp-tube-contraction}, \SeriesRef{inv:B1.1}}

For an admissible $(3,4)$ member with
$n=|12l_0-4l_1-3l_2|>1$, let $\pi:Y\to X$ be the universal
cover and $g=f\pi$. It is a finite \'etale holomorphic map of degree $n$.
Every fibre of $g$ is connected; its finite reductions are bielliptic,
its multiplicities are $3,4$, and its cusp $W_n$ normalizes to a
disjoint union of $n$ copies of $\mathrm{dP}_6$.
For $\mathcal E_X$ of Definition~\ref{I:def:coarse-log-cotangent}
and the identical quotient $\mathcal E_Y$ on $Y$, pullback induces isomorphisms
\[
R^qf_*\OO_X\simeq R^qg_*\OO_Y,\qquad
R^qf_*\mathcal E_X\simeq R^qg_*\mathcal E_Y.
\]
\end{proposition}


\begin{proof}
\begin{enumerate}[label=\arabic*.,leftmargin=*]
\item We first derive the fundamental-group presentation from the fillings.
The smooth family over the punctured base has zero section, so its
fundamental group is $\Lambda\rtimes\langle x,y\rangle$, where
$x,y$ are the analytic peripheral generators fixed by the root-deck convention and act by $A_1,A_2$.
On a finite filling the lifted affine generator has $m_j$th power
translation by $v_j$: its $m_j$ translations add to $v_j$ because
$A_jv_j=v_j$. Thus filling the normal disc imposes
$x^3=t_{v_1}$ and $y^4=t_{v_2}$, with the signs fixed by the
clockwise convention on the first page.

\item The toric covering space of the cusp tube is simply connected.
Indeed its dense torus has cocharacter lattice $\Z^3$; a loop may
be perturbed into that torus, and the normal discs of the three rays
$(0,1),(e_1,1),(e_2,1)$ kill a basis of this lattice.
This argument applies also to the infinite fan, since a compact loop
and a homotopy meet only finitely many of its locally finite charts.
The same conclusion holds for the small tube: if a loop in
$|t_c|<\epsilon$ bounds a disc in the toric space, the torus action
$(x_1,x_2,t_c)\mapsto(x_1,x_2,a t_c)$ for a sufficiently small
positive $a$ moves that compact disc into the tube, and the scaling
homotopy of its boundary stays in the tube.
The cusp quotient therefore has fundamental group
$\Lambda/\Lambda_{\rm tor}$, and its toric meridian bounds the
normal disc at fixed $x_1,x_2$. A translated collar makes this
meridian $t_\mu y^{-1}x^{-1}$, where $\gamma(\mu)=l_0$.
Van Kampen consequently kills $\Lambda_{\rm tor}$ and imposes
$xy=t_\mu$.

\item The normal closure of $\Lambda_{\rm tor}$ in the fibre lattice is
$\ker\gamma$: it contains $\hat w,\hat\delta$ and
$A_1\hat w+\hat w=\hat u$, and all these relations lie in
$\ker\gamma$. The remaining generator $c$, the image of
$\hat\gamma$, is central because $\gamma A_j=\gamma$.
The presentation is thus
$\langle c,x,y\mid c\text{ central},\ xy=c^{l_0},\
x^3=c^{l_1},\ y^4=c^{l_2}\rangle$.
Eliminating $y$ makes it abelian, with relation matrix
\[
 \begin{pmatrix}3&-l_1\\4&l_2-4l_0\end{pmatrix},
 \qquad\det=-p.
\]
Its first Smith divisor is one because its entries include three
and four; hence the group is $\Z/|p|$.
After imposing $c=1$ the quotient has $x^3=x^4=1$, so it is trivial.
Thus the image of the regular-fibre lattice, containing $c$, generates
the entire group $\Z/n$. The same is true of each reduced special
fibre. At a finite point the cyclic model is a smooth torus family
over a disc, with a free affine finite-group action. An equivariant
smooth trivialization, obtained by averaging a horizontal connection,
contracts the disc to its centre and descends to a deformation
retraction of the local total space onto $S_j$. Thus the image of
the regular-fibre lattice in the local fundamental group is contained
in the image of $\pi_1(S_j)$; the same holds after mapping to
$\pi_1(X)$. At the cusp $\hat\gamma$ survives in
$\Lambda/\Lambda_{\rm tor}$. The explicit normal-crossing contraction
of the tube onto $W$ in Lemma~\ref{I:lem:cusp-tube-contraction} identifies their fundamental groups,
so the same generator is in the image of $\pi_1(W)$.
The covering-space criterion therefore makes the inverse image of
every fibre connected. Admissibility gives $3\nmid l_1$ and odd $l_2$,
so $\gcd(n,12)=1$. The etale cover preserves fibre multiplicities.
On a finite reduction it cannot kill the nontrivial order-$3$ or
order-$4$ linear holonomy, since the covering degree is prime to it;
the resulting finite torus quotient is again bielliptic.

\item Normalization commutes with the finite etale base change $Y\to X$.
The pullback of the simply connected normalization $\mathrm{dP}_6$
is its trivial degree-$n$ cover. The double curves and triple points
also pull back etale, and their incidence maps are the lifts of the
original incidence maps. The lifted incidence maps therefore make the dual complex a
degree-$n$ covering of the original real two-torus. It is connected because
$W_n$ is connected. The normalization resolution for $\OO_{W_n}$
has acyclic components, double curves and point strata. Its global
cochain complex is consequently the cellular cochain complex of this
dual torus, with $n,3n,2n$ cells. Thus its coherent cohomology is
$(1,2,1)$, and pullback from $W$ is an isomorphism over $\C$:
on a finite cover of real tori it is induced by a full-rank map of
their integral first cohomology lattices.

\item Modulo $d\log q$, the two toric forms $d\log x_i$ are invariant
under all cusp gluings as computed in \ref{inv:B1.1}. Hence
$\mathcal E_X|_W=\OO_W^2$ and $\mathcal E_Y|_{W_n}=\OO_{W_n}^2$,
and the same pullback argument applies to their cohomology.
Over a regular fibre, an isogeny of complex tori induces isomorphisms
on $H^q(\OO)$ and $H^q(\Omega^1)$ over $\C$.
At a finite fibre, make the cyclic base change. The map is an isogeny
of the smooth torus families and induces isomorphisms of their
cohomology bundles. Taking the same cyclic invariant eigenspaces
descends these isomorphisms, including all nilpotent layers of the
scheme fibre. Grauert base change at the cusp, and this explicit
calculation at the finite points, prove the asserted bundle maps are
isomorphisms on all of $B$.
\end{enumerate}
\end{proof}

% inv:B3.3 ORIGINAL
\begin{proposition}[Canonical rings and rational topology of the covers]\hypertarget{tag-01S9}{}\label{inv:B3.3}
\StatementRecord{inv:B3.3}{proved}{01S9}{\SeriesRef{inv:B3.1}, \SeriesRef{inv:B2.2}, \SeriesRef{inv:B3.2}}

For $X$ and $Y$ as above,
\[
K_Y=\pi^*K_X=\OO_Y(-2S'_2),\quad
\bigoplus_{m\ge0}H^0(Y,-mK_Y)=\C[U,V],\quad
\deg(U,V)=(1,2).
\]
The same ring holds for $X$.
Over $\C$, the Betti numbers of $Y$ are
$(1,0,n-1,0,n-1,0,1)$; its Fr\"olicher sequence degenerates at $E_2$.
For every nontrivial character of the deck group, the character
part of Dolbeault cohomology consists of one class in $(1,1)$ and
one in $(2,2)$ and no others.
\end{proposition}


\begin{proof}
\begin{enumerate}[label={$\langle1\rangle$\arabic*.},leftmargin=*]
\item The canonical formula pulls back through an etale map, and
$4S'_2=g^*p_2$. \ref{inv:B3.1} gives $g_*\OO_Y=\OO_B$.
The valuation argument in \ref{inv:B2.2} therefore applies unchanged and proves
the ring statement with all multiplication maps.

\item The comparisons of \ref{inv:B3.1} are pullback maps and identify the logarithmic
cohomology bundles with the invariant part; the deck group acts
trivially on those bundles. At the cusp its action on the $n$
normalization components is the regular permutation representation.
The quotient by the diagonal has each nontrivial character once.
Thus the first-row additional skyscraper in \ref{inv:B3.2} has precisely these
characters, and duality gives the conjugate characters in its
second-row additional skyscraper. All other cohomology is invariant
and equal to that of $X$. Coherent Leray now proves the assertion
about the Dolbeault character spaces.

\item The de Rham complex and its Fr\"olicher sequence split under the
idempotents of the finite deck group. A nontrivial character has just
the two stated Dolbeault entries. No differential can connect them
or connect either to another nonzero entry, so both survive.
The invariant summand is the de Rham complex of $X$ by finite
etale descent and averaging. We compute its rational topology here
from the actual attachment maps; the integral homology assertion
of AC26 is not used as an input. In the remainder of this proof
all homology and cohomology coefficients are $\Q$.

\item Let $N$ be a small closed cusp tube, $M=\partial N$, and
$U=X\setminus\operatorname{int}N$. The space $U$ is covered by
two pieces retracting onto $S_1,S_2$, whose intersection retracts
onto a regular fibre $F$. The retractions are obtained from the
equivariant smooth torus trivializations used in \ref{inv:B3.1}. The maps
$F\to S_j$ are the central finite covering maps; translation of
their origins induces the identity on their rational cohomology.
Finite-cover transfer identifies their pullbacks with the invariant
subspaces for $A_j$. Put
\[
 \psi_1=2u+w+3\delta,\qquad \psi_2=u+w+2\delta,
 \qquad \eta=u\wedge w+6\gamma\wedge\delta.
\]
Direct substitution in $A_1,A_2$ shows that the invariant
degree-one spaces are $\langle\gamma,\psi_j\rangle$ and the
invariant degree-two spaces are
$\langle\gamma\wedge\psi_j,\eta\rangle$.
Their dimensions are two: the nontrivial eigenvalues of $A_j$ on
the complementary plane are the inverse primitive third or fourth
roots, so the exterior-square fixed space has exactly the two
displayed generators.
In the transfer-normalized dual homology bases, the two
Mayer--Vietoris maps are therefore the following actual matrices:
\[
 \alpha_1=\begin{pmatrix}
 1&0&0&0\\0&2&1&3\\-1&0&0&0\\0&-1&-1&-2
 \end{pmatrix},\qquad
 \alpha_2=\begin{pmatrix}
 2&1&3&0&0&0\\0&0&6&1&0&0\\
 -1&-1&-2&0&0&0\\0&0&-6&-1&0&0
 \end{pmatrix}.
\]
The second matrix uses the ordered wedge basis
$(\hat\gamma\hat u,\hat\gamma\hat w,\hat\gamma\hat\delta,
\hat u\hat w,\hat u\hat\delta,\hat w\hat\delta)$.
Both matrices have rank three, and
$\ker\alpha_1=\Q\nu_0$, with
$\nu_0=-\hat u-\hat w+\hat\delta$.
Consequently
\[
 H_1(U)=\Q,\qquad
 0\longrightarrow\operatorname{coker}\alpha_2=\Q
 \longrightarrow H_2(U)
 \xrightarrow{\partial}\Q\nu_0\longrightarrow0.       \tag{B.1}
\]
The generator $\Theta\in H^1(U)$ is normalized by
$\Theta|_F=\gamma$. The finite affine generators have
$3\bar g_1=l_1c$, $4\bar g_2=l_2c$ in rational homology, where
$c$ is the image of $\hat\gamma$ and $\Theta(c)=1$.
For the toric meridian $a_0=t_\mu y^{-1}x^{-1}$ this gives the
explicit collar value
\[
 \Theta(j_*a_0)=l_0-l_1/3-l_2/4=p/12.                 \tag{B.2}
\]

\item We justify the comparison with $N$ without assuming a collapse
theorem from AC26. In each normal-crossing chart
$t_c=z_1\cdots z_r$, $1\le r\le3$, choose the real vector field
whose $z_i$ component is
$-z_i|z_i|^{-2}/\sum_j|z_j|^{-2}$ off the divisor.
It satisfies $\dot t_c=-t_c$ and has norm at most
$C|t_c|^{1/r}$. A finite partition of unity preserves the first
identity and gives a global field with norm at most
$C'|t_c|^{1/3}$ after shrinking the tube. Extend it by zero on $W$.
Its flow exists for all positive time, and its tail length is
bounded by $3C'\epsilon^{1/3}e^{-T/3}$ after time $T$.
Thus the flow converges uniformly to $W$ and gives a continuous
strong deformation retraction fixing $W$. The same estimate proves
continuity at the divisor and at the endpoint of the retraction.

\item The constant-sheaf normalization resolution of $W$ has strata
$\mathrm{dP}_6$, three copies of $\Pone$, and two points.
Its row of degree-zero cohomology is
\[
 \Q\xrightarrow{0}\Q^3
 \xrightarrow{(a,b,c)\mapsto(a-b+c,a-b+c)}\Q^2.
\]
The strata have no degree-one cohomology. This resolution therefore
gives $H_1(N)=H_1(W)=\Q^2$, $H_4(N)=\Q[W]$ and $H_5(N)=0$;
it also gives $e(W)=6-3\cdot2+2=2$.
Poincare--Lefschetz duality for the compact smooth manifold with
boundary $(N,M)$ now shows that $H_1(M)\to H_1(N)$ and
$H_2(M)\to H_2(N)$ are surjective. For degree one the next
relative group is $H_1(N,M)=H^5(N)=0$.
For degree two the map to $H_2(N,M)=H^4(N)$ is the intersection
pairing with $[W]$, which is zero because $W$ is the divisor of
the global function $t_c$ on $N$. The Wang sequence for the
boundary mapping torus, with $\rank(M_0-I)=2$, gives
$\dim H_1(M)=3$. Its kernel to $H_1(N)$ is exactly
$\Q a_0$: this meridian is nonzero under projection to the base
circle and bounds its normal disc in $N$.

\item We exhibit two classes in
$K_2=\ker(H_2(M)\to H_2(N))$.
For $\nu=\hat w,\hat\delta$, let $\xi_\nu$ be the two-torus
obtained by translating the toric meridian around the angular
circle $S^1_\nu$. In the chart $t_c=z_0z_1z_2$ corresponding to
the rays $(0,1),(e_1,1),(e_2,1)$, these are the boundaries of
$D^2_{z_0}\times S^1_{z_i}$ with the other coordinate fixed
and the radii chosen so that $|t_c|=\epsilon$ on the boundary.
They therefore bound in $N$ with the literal toric meridian;
no choice of an unrecorded collar lift is being made.
Cutting the meridian into the two arcs of the cover defining
(B.1), the two cylinders have end circles differing by
$(I-A_1^{-1})\nu$. With this orientation of the Mayer--Vietoris
boundary, matrix multiplication gives
\[
 \partial j_*\xi_{\hat w}=-\nu_0,\qquad
 \partial j_*\xi_{\hat\delta}=0.                       \tag{B.3}
\]
The affine translation of either endpoint has no effect on its
homology class, which explains why this boundary formula is the
same for every translation twist.

\item Choose $\mathcal A\in H^2(U)$ restricting to $\eta$ on $F$;
it exists by Mayer--Vietoris because the two transfer-normalized
classes $\eta$ have equal restrictions. Fibrewise translation by
$S^1_{\hat\delta}$ defines a global smooth action on $U$:
$A_j\hat\delta=\hat\delta$, and all affine translations commute
with this circle. Slant product of its action map with the circle
defines $\iota:H^2(U)\to H^1(U)$. Naturality on $F$ gives
$\iota\mathcal A|_F=\iota_{\hat\delta}\eta=-6\gamma$.
The restriction $H^1(U)\to H^1(F)$ is injective by (B.1), so
$\iota\mathcal A=-6\Theta$. By the definition of slant product,
\[
 \langle\mathcal A,j_*\xi_{\hat\delta}\rangle
 =\langle\iota\mathcal A,j_*a_0\rangle=-p/2.           \tag{B.4}
\]
We may add the class pulled back from $(\Q\nu_0)^*$ under
the boundary map in (B.1) so that
$\langle\mathcal A,j_*\xi_{\hat w}\rangle=0$;
this does not change its restriction to $F$ or (B.4).
In the coordinates on $H_2(U)$ given by $\partial/\nu_0$ and
this normalized pairing, the two actual attachment cycles have
matrix
\[
 \bigl[j_*\xi_{\hat w}\ \ j_*\xi_{\hat\delta}\bigr]
 =\begin{pmatrix}-1&0\\0&-p/2\end{pmatrix}.             \tag{B.5}
\]
It is invertible over $\Q$ because admissibility implies
$p\ne0$. Thus the image of $K_2$ already fills $H_2(U)$.
Likewise (B.2) says the image of $\ker(H_1(M)\to H_1(N))$
is nonzero in $H_1(U)$. Using the surjectivity to $N$ just proved,
Mayer--Vietoris for $X=U\cup N$ now gives $H_1(X)=H_2(X)=0$.
Every other fibre has Euler number zero, so additivity over the
three special fibres and the smooth torus bundle gives $e(X)=2$.
Poincare duality then gives $b_4=b_5=0$ and $b_3=0$.
This proves the asserted rational homology of $X$ for all admissible
twists; the only changing determinant in (B.2)--(B.5) is a nonzero
constant multiple of $p$. The integer matrices and all exterior-power
identities were checked exactly in \texttt{code/b\_topology.py},
with output \texttt{outputs/b\_topology.txt}.

\item Part C's argument using this rational topology and the same Hodge
table now gives degeneration of the invariant summand at $E_2$.
Summing it with the nontrivial character spaces proves the Betti
numbers and the degeneration statement for $Y$. In particular
$e(Y)=2n$, as also follows from the covering degree.
\end{enumerate}
\end{proof}

% CANDIDATE11
\begin{proposition}[Canonical rings and rational topology of the covers]\hypertarget{tag-01S9}{}\label{inv:B3.3}
\StatementRecord{inv:B3.3}{proved}{01S9}{\SeriesRef{inv:B3.1}, \SeriesRef{inv:B2.2}, \SeriesRef{inv:B3.2}}

For $X$ and $Y$ as above,
\[
K_Y=\pi^*K_X=\OO_Y(-2S'_2),\quad
\bigoplus_{m\ge0}H^0(Y,-mK_Y)=\C[U,V],\quad
\deg(U,V)=(1,2).
\]
The same ring holds for $X$.
Over $\C$, the Betti numbers of $Y$ are
$(1,0,n-1,0,n-1,0,1)$; its Fr\"olicher sequence degenerates at $E_2$.
For every nontrivial character of the deck group, the character
part of Dolbeault cohomology consists of one class in $(1,1)$ and
one in $(2,2)$ and no others.
\end{proposition}


\begin{proof}
\begin{enumerate}[label=\arabic*.,leftmargin=*]
\item The canonical formula pulls back through an etale map, and
$4S'_2=g^*p_2$. \ref{inv:B3.1} gives $g_*\OO_Y=\OO_B$.
The valuation argument in \ref{inv:B2.2} therefore applies unchanged and proves
the ring statement with all multiplication maps.

\item The comparisons of \ref{inv:B3.1} are pullback maps and identify the logarithmic
cohomology bundles with the invariant part; the deck group acts
trivially on those bundles. At the cusp its action on the $n$
normalization components is the regular permutation representation.
The quotient by the diagonal has each nontrivial character once.
Thus the first-row additional skyscraper in \ref{inv:B3.2} has precisely these
characters, and duality gives the conjugate characters in its
second-row additional skyscraper. All other cohomology is invariant
and equal to that of $X$. Coherent Leray now proves the assertion
about the Dolbeault character spaces.

\item The de Rham complex and its Fr\"olicher sequence split under the
idempotents of the finite deck group. A nontrivial character has just
the two stated Dolbeault entries. No differential can connect them
or connect either to another nonzero entry, so both survive.
The invariant summand is the de Rham complex of $X$ by finite
etale descent and averaging. Its rational topology is computed by the following attachment maps. In the remainder of this proof
all homology and cohomology coefficients are $\Q$.

\item Let $N$ be a small closed cusp tube, $M=\partial N$, and
$U=X\setminus\operatorname{int}N$. The space $U$ is covered by
two pieces retracting onto $S_1,S_2$, whose intersection retracts
onto a regular fibre $F$. The retractions are obtained from the
equivariant smooth torus trivializations used in \ref{inv:B3.1}. The maps
$F\to S_j$ are the central finite covering maps; translation of
their origins induces the identity on their rational cohomology.
Finite-cover transfer identifies their pullbacks with the invariant
subspaces for $A_j$. Put
\[
 \psi_1=2u+w+3\delta,\qquad \psi_2=u+w+2\delta,
 \qquad \eta=u\wedge w+6\gamma\wedge\delta.
\]
Direct substitution in $A_1,A_2$ shows that the invariant
degree-one spaces are $\langle\gamma,\psi_j\rangle$ and the
invariant degree-two spaces are
$\langle\gamma\wedge\psi_j,\eta\rangle$.
Their dimensions are two: the nontrivial eigenvalues of $A_j$ on
the complementary plane are the inverse primitive third or fourth
roots, so the exterior-square fixed space has exactly the two
displayed generators.
In the transfer-normalized dual homology bases, the two
Mayer--Vietoris maps are therefore the following matrices:
\[
 \alpha_1=\begin{pmatrix}
 1&0&0&0\\0&2&1&3\\-1&0&0&0\\0&-1&-1&-2
 \end{pmatrix},\qquad
 \alpha_2=\begin{pmatrix}
 2&1&3&0&0&0\\0&0&6&1&0&0\\
 -1&-1&-2&0&0&0\\0&0&-6&-1&0&0
 \end{pmatrix}.
\]
The second matrix uses the ordered wedge basis
$(\hat\gamma\hat u,\hat\gamma\hat w,\hat\gamma\hat\delta,
\hat u\hat w,\hat u\hat\delta,\hat w\hat\delta)$.
Both matrices have rank three, and
$\ker\alpha_1=\Q\nu_0$, with
$\nu_0=-\hat u-\hat w+\hat\delta$.
Consequently
\begin{equation}
 H_1(U)=\Q,\qquad
 0\longrightarrow\operatorname{coker}\alpha_2=\Q
 \longrightarrow H_2(U)
 \xrightarrow{\partial}\Q\nu_0\longrightarrow0.
\label{pub:I:eq:B-1}
\end{equation}
The generator $\Theta\in H^1(U)$ is normalized by
$\Theta|_F=\gamma$. The finite affine generators have
$3\bar g_1=l_1c$, $4\bar g_2=l_2c$ in rational homology, where
$c$ is the image of $\hat\gamma$ and $\Theta(c)=1$.
For the toric meridian $a_0=t_\mu y^{-1}x^{-1}$ this gives the
explicit collar value
\begin{equation}
 \Theta(j_*a_0)=l_0-l_1/3-l_2/4=p/12.
\label{pub:I:eq:B-2}
\end{equation}

\item The comparison with $N$ is given by this explicit deformation
retraction. In each normal-crossing chart
$t_c=z_1\cdots z_r$, $1\le r\le3$, choose the real vector field
whose $z_i$ component is
$-z_i|z_i|^{-2}/\sum_j|z_j|^{-2}$ off the divisor.
It satisfies $\dot t_c=-t_c$ and has norm at most
$C|t_c|^{1/r}$. A finite partition of unity preserves the first
identity and gives a global field with norm at most
$C'|t_c|^{1/3}$ after shrinking the tube. Extend it by zero on $W$.
Its flow exists for all positive time, and its tail length is
bounded by $3C'\epsilon^{1/3}e^{-T/3}$ after time $T$.
Thus the flow converges uniformly to $W$ and gives a continuous
strong deformation retraction fixing $W$. The same estimate proves
continuity at the divisor and at the endpoint of the retraction.

\item The constant-sheaf normalization resolution of $W$ has strata
$\mathrm{dP}_6$, three copies of $\Pone$, and two points.
Its row of degree-zero cohomology is
\[
 \Q\xrightarrow{0}\Q^3
 \xrightarrow{(a,b,c)\mapsto(a-b+c,a-b+c)}\Q^2.
\]
The strata have no degree-one cohomology. This resolution therefore
gives $H_1(N)=H_1(W)=\Q^2$, $H_4(N)=\Q[W]$ and $H_5(N)=0$;
it also gives $e(W)=6-3\cdot2+2=2$.
Poincare--Lefschetz duality for the compact smooth manifold with
boundary $(N,M)$ now shows that $H_1(M)\to H_1(N)$ and
$H_2(M)\to H_2(N)$ are surjective. For degree one the next
relative group is $H_1(N,M)=H^5(N)=0$.
For degree two the map to $H_2(N,M)=H^4(N)$ is the intersection
pairing with $[W]$, which is zero because $W$ is the divisor of
the global function $t_c$ on $N$. The Wang sequence for the
boundary mapping torus, with $\rank(M_0-I)=2$, gives
$\dim H_1(M)=3$. Its kernel to $H_1(N)$ is exactly
$\Q a_0$: this meridian is nonzero under projection to the base
circle and bounds its normal disc in $N$.

\item We exhibit two classes in
$K_2=\ker(H_2(M)\to H_2(N))$.
For $\nu=\hat w,\hat\delta$, let $\xi_\nu$ be the two-torus
obtained by translating the toric meridian around the angular
circle $S^1_\nu$. In the chart $t_c=z_0z_1z_2$ corresponding to
the rays $(0,1),(e_1,1),(e_2,1)$, these are the boundaries of
$D^2_{z_0}\times S^1_{z_i}$ with the other coordinate fixed
and the radii chosen so that $|t_c|=\epsilon$ on the boundary.
They therefore bound in $N$ with the specified toric meridian.
Cutting the meridian into the two arcs of the cover defining
\eqref{pub:I:eq:B-1}, the two cylinders have end circles differing by
$(I-A_1^{-1})\nu$. With this orientation of the Mayer--Vietoris
boundary, matrix multiplication gives
\begin{equation}
 \partial j_*\xi_{\hat w}=-\nu_0,\qquad
 \partial j_*\xi_{\hat\delta}=0.
\label{pub:I:eq:B-3}
\end{equation}
The affine translation of either endpoint has no effect on its
homology class, which explains why this boundary formula is the
same for every translation twist.

\item Choose $\mathcal A\in H^2(U)$ restricting to $\eta$ on $F$;
it exists by Mayer--Vietoris because the two transfer-normalized
classes $\eta$ have equal restrictions. Fibrewise translation by
$S^1_{\hat\delta}$ defines a global smooth action on $U$:
$A_j\hat\delta=\hat\delta$, and all affine translations commute
with this circle. Slant product of its action map with the circle
defines $\iota:H^2(U)\to H^1(U)$. Naturality on $F$ gives
$\iota\mathcal A|_F=\iota_{\hat\delta}\eta=-6\gamma$.
The restriction $H^1(U)\to H^1(F)$ is injective by \eqref{pub:I:eq:B-1}, so
$\iota\mathcal A=-6\Theta$. By the definition of slant product,
\begin{equation}
 \langle\mathcal A,j_*\xi_{\hat\delta}\rangle
 =\langle\iota\mathcal A,j_*a_0\rangle=-p/2.
\label{pub:I:eq:B-4}
\end{equation}
We may add the class pulled back from $(\Q\nu_0)^*$ under
the boundary map in \eqref{pub:I:eq:B-1} so that
$\langle\mathcal A,j_*\xi_{\hat w}\rangle=0$;
this does not change its restriction to $F$ or \eqref{pub:I:eq:B-4}.
In the coordinates on $H_2(U)$ given by $\partial/\nu_0$ and
this normalized pairing, the two attachment cycles have
matrix
\begin{equation}
 \bigl[j_*\xi_{\hat w}\ \ j_*\xi_{\hat\delta}\bigr]
 =\begin{pmatrix}-1&0\\0&-p/2\end{pmatrix}.
\label{pub:I:eq:B-5}
\end{equation}
It is invertible over $\Q$ because admissibility implies
$p\ne0$. Thus the image of $K_2$ already fills $H_2(U)$.
Likewise \eqref{pub:I:eq:B-2} says the image of $\ker(H_1(M)\to H_1(N))$
is nonzero in $H_1(U)$. Using the surjectivity to $N$ just proved,
Mayer--Vietoris for $X=U\cup N$ now gives $H_1(X)=H_2(X)=0$.
Every other fibre has Euler number zero, so additivity over the
three special fibres and the smooth torus bundle gives $e(X)=2$.
Poincare duality then gives $b_4=b_5=0$ and $b_3=0$.
This proves the asserted rational homology of $X$ for all admissible
twists; the only changing determinant in \eqref{pub:I:eq:B-2}--\eqref{pub:I:eq:B-5} is a nonzero
constant multiple of $p$. 

\item Proposition~\ref{inv:C1.1}, applied to this rational topology and
the same Hodge table, gives degeneration of the invariant summand at $E_2$.
Summing it with the nontrivial character spaces proves the Betti
numbers and the degeneration statement for $Y$. In particular
$e(Y)=2n$, as also follows from the covering degree.
\end{enumerate}
\end{proof}

% PROPOSED
\begin{proposition}[Canonical rings and rational topology of the covers]\hypertarget{tag-01S9}{}\label{inv:B3.3}
\StatementRecord{inv:B3.3}{proved}{01S9}{\SeriesRef{inv:B3.1}, \SeriesRef{inv:B2.2}, \SeriesRef{inv:B3.2}}

For $X$ and $Y$ as above,
\[
K_Y=\pi^*K_X=\OO_Y(-2S'_2),\quad
\bigoplus_{m\ge0}H^0(Y,-mK_Y)=\C[U,V],\quad
\deg(U,V)=(1,2).
\]
The same ring holds for $X$.
Over $\C$, the Betti numbers of $Y$ are
$(1,0,n-1,0,n-1,0,1)$; its Fr\"olicher sequence degenerates at $E_2$.
For every nontrivial character of the deck group, the character
part of Dolbeault cohomology consists of one class in $(1,1)$ and
one in $(2,2)$ and no others.
\end{proposition}


\begin{proof}
\begin{enumerate}[label=\arabic*.,leftmargin=*]
\item The canonical formula pulls back through an etale map, and
$4S'_2=g^*p_2$. \ref{inv:B3.1} gives $g_*\OO_Y=\OO_B$.
The valuation argument in \ref{inv:B2.2} therefore applies unchanged and proves
the ring statement with all multiplication maps.

\item The comparisons of \ref{inv:B3.1} are pullback maps and identify the logarithmic
cohomology bundles with the invariant part; the deck group acts
trivially on those bundles. At the cusp its action on the $n$
normalization components is the regular permutation representation.
The quotient by the diagonal has each nontrivial character once.
Thus the first-row additional skyscraper in \ref{inv:B3.2} has precisely these
characters, and duality gives the conjugate characters in its
second-row additional skyscraper. All other cohomology is invariant
and equal to that of $X$. Coherent Leray now proves the assertion
about the Dolbeault character spaces.

\item The de Rham complex and its Fr\"olicher sequence split under the
idempotents of the finite deck group. A nontrivial character has just
the two stated Dolbeault entries. No differential can connect them
or connect either to another nonzero entry, so both survive.
The invariant summand is the de Rham complex of $X$ by finite
etale descent and averaging. Its rational topology is computed by the following attachment maps. In the remainder of this proof
all homology and cohomology coefficients are $\Q$.

\item Let $N$ be a small closed cusp tube, $M=\partial N$, and
$U=X\setminus\operatorname{int}N$. The space $U$ is covered by
two pieces retracting onto $S_1,S_2$, whose intersection retracts
onto a regular fibre $F$. The retractions are obtained from the
equivariant smooth torus trivializations used in \ref{inv:B3.1}. The maps
$F\to S_j$ are the central finite covering maps; translation of
their origins induces the identity on their rational cohomology.
Finite-cover transfer identifies their pullbacks with the invariant
subspaces for $A_j$. Put
\[
 \psi_1=2u+w+3\delta,\qquad \psi_2=u+w+2\delta,
 \qquad \eta=u\wedge w+6\gamma\wedge\delta.
\]
Direct substitution in $A_1,A_2$ shows that the invariant
degree-one spaces are $\langle\gamma,\psi_j\rangle$ and the
invariant degree-two spaces are
$\langle\gamma\wedge\psi_j,\eta\rangle$.
Their dimensions are two: the nontrivial eigenvalues of $A_j$ on
the complementary plane are the inverse primitive third or fourth
roots, so the exterior-square fixed space has exactly the two
displayed generators.
In the transfer-normalized dual homology bases, the two
Mayer--Vietoris maps are therefore the following matrices:
\[
 \alpha_1=\begin{pmatrix}
 1&0&0&0\\0&2&1&3\\-1&0&0&0\\0&-1&-1&-2
 \end{pmatrix},\qquad
 \alpha_2=\begin{pmatrix}
 2&1&3&0&0&0\\0&0&6&1&0&0\\
 -1&-1&-2&0&0&0\\0&0&-6&-1&0&0
 \end{pmatrix}.
\]
The second matrix uses the ordered wedge basis
$(\hat\gamma\hat u,\hat\gamma\hat w,\hat\gamma\hat\delta,
\hat u\hat w,\hat u\hat\delta,\hat w\hat\delta)$.
Both matrices have rank three, and
$\ker\alpha_1=\Q\nu_0$, with
$\nu_0=-\hat u-\hat w+\hat\delta$.
Consequently
\begin{equation}
 H_1(U)=\Q,\qquad
 0\longrightarrow\operatorname{coker}\alpha_2=\Q
 \longrightarrow H_2(U)
 \xrightarrow{\partial}\Q\nu_0\longrightarrow0.
\label{pub:I:eq:B-1}
\end{equation}
The generator $\Theta\in H^1(U)$ is normalized by
$\Theta|_F=\gamma$. The finite affine generators have
$3\bar g_1=l_1c$, $4\bar g_2=l_2c$ in rational homology, where
$c$ is the image of $\hat\gamma$ and $\Theta(c)=1$.
For the toric meridian $a_0=t_\mu y^{-1}x^{-1}$ this gives the
explicit collar value
\begin{equation}
 \Theta(j_*a_0)=l_0-l_1/3-l_2/4=p/12.
\label{pub:I:eq:B-2}
\end{equation}

\item Lemma~\ref{I:lem:cusp-tube-contraction} gives the strong
deformation retraction $N\to W$, fixing $W$ pointwise.

\item The constant-sheaf normalization resolution of $W$ has strata
$\mathrm{dP}_6$, three copies of $\Pone$, and two points.
Its row of degree-zero cohomology is
\[
 \Q\xrightarrow{0}\Q^3
 \xrightarrow{(a,b,c)\mapsto(a-b+c,a-b+c)}\Q^2.
\]
The strata have no degree-one cohomology. This resolution therefore
gives $H_1(N)=H_1(W)=\Q^2$, $H_4(N)=\Q[W]$ and $H_5(N)=0$;
it also gives $e(W)=6-3\cdot2+2=2$.
Poincare--Lefschetz duality for the compact smooth manifold with
boundary $(N,M)$ now shows that $H_1(M)\to H_1(N)$ and
$H_2(M)\to H_2(N)$ are surjective. For degree one the next
relative group is $H_1(N,M)=H^5(N)=0$.
For degree two the map to $H_2(N,M)=H^4(N)$ is the intersection
pairing with $[W]$, which is zero because $W$ is the divisor of
the global function $t_c$ on $N$. The Wang sequence for the
boundary mapping torus, with $\rank(M_0-I)=2$, gives
$\dim H_1(M)=3$. Its kernel to $H_1(N)$ is exactly
$\Q a_0$: this meridian is nonzero under projection to the base
circle and bounds its normal disc in $N$.

\item We exhibit two classes in
$K_2=\ker(H_2(M)\to H_2(N))$.
For $\nu=\hat w,\hat\delta$, let $\xi_\nu$ be the two-torus
obtained by translating the toric meridian around the angular
circle $S^1_\nu$. In the chart $t_c=z_0z_1z_2$ corresponding to
the rays $(0,1),(e_1,1),(e_2,1)$, these are the boundaries of
$D^2_{z_0}\times S^1_{z_i}$ with the other coordinate fixed
and the radii chosen so that $|t_c|=\epsilon$ on the boundary.
They therefore bound in $N$ with the specified toric meridian.
Cutting the meridian into the two arcs of the cover defining
\eqref{pub:I:eq:B-1}, the two cylinders have end circles differing by
$(I-A_1^{-1})\nu$. With this orientation of the Mayer--Vietoris
boundary, matrix multiplication gives
\begin{equation}
 \partial j_*\xi_{\hat w}=-\nu_0,\qquad
 \partial j_*\xi_{\hat\delta}=0.
\label{pub:I:eq:B-3}
\end{equation}
The affine translation of either endpoint has no effect on its
homology class, which explains why this boundary formula is the
same for every translation twist.

\item Choose $\mathcal A\in H^2(U)$ restricting to $\eta$ on $F$;
it exists by Mayer--Vietoris because the two transfer-normalized
classes $\eta$ have equal restrictions. Fibrewise translation by
$S^1_{\hat\delta}$ defines a global smooth action on $U$:
$A_j\hat\delta=\hat\delta$, and all affine translations commute
with this circle. Slant product of its action map with the circle
defines $\iota:H^2(U)\to H^1(U)$. Naturality on $F$ gives
$\iota\mathcal A|_F=\iota_{\hat\delta}\eta=-6\gamma$.
The restriction $H^1(U)\to H^1(F)$ is injective by \eqref{pub:I:eq:B-1}, so
$\iota\mathcal A=-6\Theta$. By the definition of slant product,
\begin{equation}
 \langle\mathcal A,j_*\xi_{\hat\delta}\rangle
 =\langle\iota\mathcal A,j_*a_0\rangle=-p/2.
\label{pub:I:eq:B-4}
\end{equation}
We may add the class pulled back from $(\Q\nu_0)^*$ under
the boundary map in \eqref{pub:I:eq:B-1} so that
$\langle\mathcal A,j_*\xi_{\hat w}\rangle=0$;
this does not change its restriction to $F$ or \eqref{pub:I:eq:B-4}.
In the coordinates on $H_2(U)$ given by $\partial/\nu_0$ and
this normalized pairing, the two attachment cycles have
matrix
\begin{equation}
 \bigl[j_*\xi_{\hat w}\ \ j_*\xi_{\hat\delta}\bigr]
 =\begin{pmatrix}-1&0\\0&-p/2\end{pmatrix}.
\label{pub:I:eq:B-5}
\end{equation}
It is invertible over $\Q$ because admissibility implies
$p\ne0$. Thus the image of $K_2$ already fills $H_2(U)$.
Likewise \eqref{pub:I:eq:B-2} says the image of $\ker(H_1(M)\to H_1(N))$
is nonzero in $H_1(U)$. Using the surjectivity to $N$ just proved,
Mayer--Vietoris for $X=U\cup N$ now gives $H_1(X)=H_2(X)=0$.
Every other fibre has Euler number zero, so additivity over the
three special fibres and the smooth torus bundle gives $e(X)=2$.
Poincare duality then gives $b_4=b_5=0$ and $b_3=0$.
This proves the asserted rational homology of $X$ for all admissible
twists; the only changing determinant in \eqref{pub:I:eq:B-2}--\eqref{pub:I:eq:B-5} is a nonzero
constant multiple of $p$. 

\item Proposition~\ref{inv:C1.1}, applied to this rational topology and
the same Hodge table, gives degeneration of the invariant summand at $E_2$.
Summing it with the nontrivial character spaces proves the Betti
numbers and the degeneration statement for $Y$. In particular
$e(Y)=2n$, as also follows from the covering degree.
\end{enumerate}
\end{proof}

% EXTRACTED LEMMA
\begin{lemma}[Contraction of the cusp tube]
\hypertarget{tag-TOPOLOGY-CUSP-TAG}{}\label{I:lem:cusp-tube-contraction}
\StatementRecord{I:lem:cusp-tube-contraction}{proved}{TOPOLOGY-CUSP-TAG}{The normal-crossing charts of the proper local cusp filling}
The reduced one-component hexagonal cusp fibre $W$ in the proper local
filling has a closed tube $N=\{|t_c|\le\epsilon\}$ that strongly
deformation retracts onto $W$. The retraction fixes $W$ pointwise.
\end{lemma}

\begin{proof}
Choose $\epsilon>0$ small enough that the compact tube is covered by
finitely many normal-crossing charts
$t_c=z_1\cdots z_r$, $1\le r\le3$. Off the divisor, take the real
vector field whose $z_i$ component is
\[
 -\frac{z_i|z_i|^{-2}}{\sum_{j=1}^r|z_j|^{-2}},
\]
with its other coordinate components zero. It satisfies
$\dot t_c=-t_c$. Each component has absolute value at most
$(\sum_j|z_j|^{-2})^{-1/2}\le\min_j|z_j|
\le|t_c|^{1/r}$, so the norm is at most $C|t_c|^{1/r}$.
A finite partition of unity preserves $\dot t_c=-t_c$ and gives a
global field of norm at most $C'|t_c|^{1/3}$ after shrinking the tube
so that $|t_c|\le1$. Extend the field by zero on $W$.

Its flow $\Phi_T$ on $N\setminus W$ satisfies
$t_c(\Phi_T(x))=e^{-T}t_c(x)$. It therefore remains in the compact
tube and exists for every $T\ge0$. Its length after time $T$ is at most
\[
 3C'|t_c(x)|^{1/3}e^{-T/3}
 \le3C'\epsilon^{1/3}e^{-T/3}.
\]
The flow converges uniformly to a continuous map $r:N\to W$,
with $r|_W=\mathrm{id}_W$. The same estimate gives continuity at
$W$ after extending every $\Phi_T$ by the identity there. The map
\[
 H(x,u)=
 \begin{cases}
 \Phi_{-\log(1-u)}(x),&0\le u<1,\\
 r(x),&u=1
 \end{cases}
\]
is continuous at $u=1$ by the uniform tail estimate. It has
$H(x,0)=x$, $H(x,1)\in W$, and $H(w,u)=w$ for $w\in W$,
which proves the strong deformation retraction.
\end{proof}

